
\subsubsection{6.1 Key Findings}

\textbf{Finding 1: Venue-string classification is not a reliable method for interdisciplinary alignment corpora indexed by S2AG.}

A bibliometric study that classifies only 12\% of its corpus papers cannot construct a valid $2\times 2$ contingency table. The failure occurs silently: the classification function returns \texttt{None} for unmatched papers, and downstream analysis omits them without warning. Researchers using S2AG data for interdisciplinary corpus classification must either adopt FoS-primary classification or normalize venue strings to full proceedings name variants. This finding applies to any interdisciplinary corpus indexed by S2AG, not exclusively to the huashen218 corpus.

\textbf{Finding 2: The directional citation pattern is present in the proxy corpus.}

A ratio of 0.107 at N = 9 cross-group edges cannot be reported as a statistically confirmed finding. The asymmetry predicted by H-CitAsym is directionally present in the raw edge counts of the 33-paper proxy, and no falsifier was triggered. The direction has not been refuted.

\textbf{Finding 3: Public reading lists are not adequate bibliometric corpus proxies.}

The 12.5\% yield from GitHub-linked papers to S2AG-resolved papers is a practical finding for researchers in adjacent areas. Programmatic S2AG search via citation-of-citation traversal, rather than reading list parsing, is required to reconstruct systematic review corpora at bibliometric scale.

\subsubsection{6.2 Limitations}

\textbf{Limitation 1: All statistical predictions (P1, P2, P3) are INCONCLUSIVE --- corpus too small for statistical testing.}

Coverage = 67.3\% and cross-group edges = 9 both fall below pre-registered thresholds. The designated next step (h-e1-v2 in the research log) specifies corpus augmentation via S2AG \texttt{/paper/arXiv:2406.09264/citations} to discover papers citing Shen et al. [2024] and expand the corpus programmatically to a target of 200+ papers.

\textbf{Limitation 2: Study operates on a 33-paper proxy corpus, not the intended ~400-paper systematic review corpus.}

The full corpus underlying Shen et al. [2024] was assembled from a structured systematic review process and is not publicly available as a machine-readable paper ID list. The GitHub repository is a curated reading guide, not a structured bibliometric database. The corpus proxy mismatch is itself a documented finding (Section 5.4).

\textbf{Limitation 3: Sensitivity analysis across three schemes is not executable at Scheme 1/2 coverage of 12\%.}

The pre-registered robustness check requires classification by at least two of three schemes. Scheme 1 and Scheme 2 classify only 4/33 papers in the current implementation. The fix --- extending venue substring sets to include full proceedings name variants --- is documented in \texttt{config.py} but was not applied during h-e1 execution.

\textbf{Limitation 4: Corpus selection bias limits generalizability.}

The huashen218 corpus is a curated alignment reading list, not a random sample of ML or HCI papers. Findings are scoped to papers within this corpus. Field-level generalization --- whether the observed asymmetry extends to the broader ML/NLP and HCI literatures --- requires a matched baseline comparison (Phase 5, deferred by pipeline configuration).

\textbf{Limitation 5: Temporal dimension not verified.}

Pre- and post-2022 cohort analysis is not meaningful at N = 33. The \texttt{year} field is already retrieved in batch resolution and is available for temporal stratification once the full corpus is reconstructed.

\subsubsection{6.3 Broader Impact}

The validated S2AG bibliometric pipeline and FoS-primary classification method are applicable to any interdisciplinary corpus indexed by S2AG. All code and cached API responses are provided as supplementary material to support zero-API re-execution. The ratio 0.107 should not be cited as confirmed evidence of asymmetric citation behavior without explicit acknowledgment of the corpus scale limitation. There is no foreseeable negative impact from measuring citation behavior in a public academic corpus.

